% Part: model-theory
% Chapter: interpolation
% Section: introduction

\documentclass[../../../include/open-logic-section]{subfiles}

\begin{document}

\olfileid{mod}{int}{int}

\olsection{ভূমিকা}

ইন্টারপোলেশন উপপাদ্যটি হলো এই ফলাফল: ধরা যাক
$\Entails !A \lif !B$। তবে এমন একটি !!a{sentence} $!C$ আছে যাতে
$\Entails !A \lif !C$ এবং $\Entails !C \lif !B$। আরও বলা যায়, $!C$-তে
থাকা প্রতিটি !!{constant}, !!{function} এবং !!{predicate} ($\eq$ ছাড়া)
উভয় $!A$ ও~$!B$-তেই ঘটে। এই !!{sentence} $!C$-কে $!A$ ও~$!B$-এর
\emph{ইন্টারপোল্যান্ট} বলা হয়।

ইন্টারপোলেশন উপপাদ্য নিজস্ব কারণেই আকর্ষণীয়, কিন্তু এর প্রধান গুরুত্ব
এই যে এটি কোনো তত্ত্বে সংজ্ঞায়ন এবং দুটি সঙ্গত তত্ত্ব একত্র করলে কখন
সঙ্গত তত্ত্ব পাওয়া যায়—এই দুই ধরনের ফলাফল প্রমাণে ব্যবহৃত হয়। প্রথম
ফলটি বেথের সংজ্ঞায়ন উপপাদ্য নামে পরিচিত; দ্বিতীয়টি রবিনসনের যৌথ
সঙ্গতি উপপাদ্য।

\end{document}
